\documentclass[11pt,a4paper]{article}
\usepackage[T1]{fontenc}
\usepackage[utf8]{inputenc}
\usepackage{lmodern}
\usepackage{amsmath,amssymb,amsthm}
\usepackage[margin=1.15in]{geometry}
\usepackage[hidelinks]{hyperref}
\emergencystretch=3em

\newtheorem{theorem}{Theorem}[section]
\newtheorem{proposition}[theorem]{Proposition}
\newtheorem{lemma}[theorem]{Lemma}
\newtheorem{corollary}[theorem]{Corollary}
\theoremstyle{definition}
\newtheorem{definition}[theorem]{Definition}
\newtheorem{remark}[theorem]{Remark}

\newcommand{\Cd}{\mathcal{C}_{\delta}}
\newcommand{\cle}{\leq_{\mathrm{d}}}
\newcommand{\El}{\operatorname{El}}
\newcommand{\colim}{\operatorname{colim}}
\newcommand{\Fix}{\operatorname{Fix}}
\newcommand{\Aut}{\operatorname{Aut}}
\newcommand{\id}{\mathrm{id}}
\newcommand{\Set}{\mathbf{Set}}
\newcommand{\Ab}{\mathbf{Ab}}
\newcommand{\K}{\mathcal{K}}
\newcommand{\Z}{\mathbb{Z}}
\newcommand{\F}{\mathbb{F}}
\newcommand{\R}{\mathbb{R}}
\newcommand{\AKRone}{Theorem~[1]}

\title{Exact, total and filtered small categories\thanks{AI-produced. The
results are due to runs of GPT-6 Astra (OpenAI, ChatGPT Pro mode), and the text
was written by Claude Fable 5.2 (Anthropic) at effort max; see
Section~\ref{sec:provenance} for who proved, wrote and checked what. No
human has checked this text. To the extent that any rights exist in it, they
are dedicated to the public domain under CC0~1.0.}}
\author{}
\date{25 November 2026}

\begin{document}
\maketitle

\begin{abstract}
For a non-empty connected small category we compare four properties: being
filtered; exactness of colimits of abelian groups; Grothendieck's total
asphericity, which says that colimits of simplicial sets over the category
are homotopy colimits; and being quasi-filtered, which means that two arrows
with a common source can be completed to a commutative square and that two
parallel arrows are equalized by an integral combination of arrows whose
coefficients sum to $1$. Each implies the next. We give a category of finite
ordered hypergraphs that is totally aspherical and whose colimits are exact
over no non-zero ring, and we show that the category of finite ordered
two-graphs (sets of triples that meet every four-element set in an even
number of triples) is quasi-filtered and not totally aspherical. With the
example of Isbell and Mitchell, the three implications are strict. A
conjecture in Grothendieck's \emph{Les D\'erivateurs}, that totally
aspherical categories are filtered, is therefore false; it holds for sifted
categories, those over which colimits of sets commute with finite products.
\end{abstract}

% ---- begin sec1_intro.tex (inlined for the deposit) ----
\section{Introduction}

Let $A$ be a small category, non-empty and connected. Four properties of $A$
are compared in this paper; the definitions are in
Section~\ref{sec:classes}.
\begin{itemize}
\item $A$ is \emph{filtered}.
\item $A$ is \emph{exact}: the colimit functor from $A$-indexed diagrams of
 abelian groups to abelian groups is exact.
\item $A$ is \emph{total}: every category that receives a final functor from
 $A$ is weakly contractible. This is Grothendieck's \emph{asph\'ericit\'e
 totale} \cite{Der}; it holds if and only if every connected component of the
 category of elements of every functor $A\to\Set$ is weakly contractible
 (Proposition~\ref{prop:total}), and if and only if colimits over $A$ of
 diagrams of simplicial sets are homotopy colimits (Remark~\ref{rem:hocolim}).
\item $A$ is \emph{quasi-filtered}: pairs of arrows with a common source can be
 completed to commutative squares, and parallel pairs have affine equalizers.
\end{itemize}
Each property implies the next (Propositions~\ref{prop:incl}
and~\ref{prop:qf}). Isbell and Mitchell showed in 1973 that the first
implication cannot be reversed. In the manuscript \emph{Les D\'erivateurs}
Grothendieck asks whether total categories are filtered. The results of this
paper are the following.

\begin{theorem}\label{thm:A}
Let $\Cd$ be the category of non-empty finite linearly ordered $3$-uniform
hypergraphs
in which every non-empty set $U$ of vertices satisfies $2|U|-3e(U)>0$, where
$e(U)$ is the number of hyperedges contained in $U$, with the closed
embeddings as arrows \textup{(}Definition~\ref{def:cdelta}\textup{)}. Then $\Cd$ is
total, and for every non-zero ring $R$ the colimit functor on $\Cd$-indexed
diagrams of $R$-modules is not exact.
\end{theorem}

\begin{theorem}\label{thm:B}
The category $\mathcal T$ of finite ordered two-graphs and their embeddings
is quasi-filtered and is not total.
\end{theorem}

\begin{theorem}\label{thm:C}
Among non-empty connected small categories,
\[
 \{\text{filtered}\}\subsetneq\{\text{exact}\}\subsetneq\{\text{total}\}
 \subsetneq\{\text{quasi-filtered}\}.
\]
The three inclusions are strict by the category $\Delta_{\mathrm{inj}}$ of
the finite non-empty ordinals and increasing injections, by $\Cd$, and by
$\mathcal T$. Moreover a non-empty connected category is filtered if and
only if it is total and sifted.
\end{theorem}

Theorem~\ref{thm:A} is Theorems~\ref{thm:total} and~\ref{thm:notexact};
Theorem~\ref{thm:B} is Theorem~\ref{thm:twographs}. For
Theorem~\ref{thm:C}: the inclusions are Propositions~\ref{prop:incl}
and~\ref{prop:qf}; $\Delta_{\mathrm{inj}}$ is exact by
\cite[Theorems~1 and~3]{IM73} (their category is that of the order-preserving
injections of finite totally ordered sets; $\Delta_{\mathrm{inj}}$ is a full
cosieve of it, and common successors and affine equalizers pass to a
non-empty cosieve), and it is not filtered, because no arrow
equalizes the two arrows from the ordinal with one element to the ordinal
with two; the last assertion is Proposition~\ref{prop:sifted}.

Theorem~\ref{thm:A} separates two weakenings of filteredness that are both
classical: exactness of colimits of modules, and the property that colimits of
spaces are homotopy colimits. The proofs of total asphericity
(Theorem~\ref{thm:A}) and of its failure (Theorem~\ref{thm:B}) follow one
pattern. Both categories consist of the finite closed substructures of a
countable homogeneous structure $M$. If every open subgroup of $\Aut(M)$ is
the pointwise stabilizer of a finite closed set, then every component of
every category of elements is weakly contractible or weakly equivalent to
the category itself (Proposition~\ref{prop:reduction}). So only the
classifying space of the category matters: it is contractible for $\Cd$, and
simply connected with non-zero second homology for $\mathcal T$. Exactness is
decided by a count: in a hypergraph of $\Cd$ any $r$ hyperedges span more
than $3r/2$ vertices, so one of them has a vertex of its own.

Section~\ref{sec:manuscript} relates the results to the passage of
\emph{Les D\'erivateurs}. This paper was produced by AI models; who proved,
wrote and checked what is recorded in Section~\ref{sec:provenance}.
% ---- end sec1_intro.tex ----
% ---- begin sec2_classes.tex (inlined for the deposit) ----
\section{The four classes}\label{sec:classes}

\subsection{Conventions}

Categories are small unless they are categories of sets, groups or modules.
For a category $A$ we write $NA$ for its nerve and $BA$ for its classifying
space. A functor is a \emph{weak equivalence} if it induces a weak homotopy
equivalence of classifying spaces, and $A$ is \emph{weakly contractible} if
$BA$ is. A natural transformation between two functors induces a homotopy
between the induced maps, and a category with an initial object is
contractible \cite{Quillen}, \cite[IV.3.2 and Example~IV.3.2.2]{Kbook}. For a
functor $u\colon A\to B$ and an object $b$ of $B$, the objects of
$b\backslash u$ are the pairs $(a,\ b\to u(a))$. The functor $u$ is
\emph{final} if every $b\backslash u$ is non-empty and connected
\cite[IX.3]{CWM}. Quillen's Theorem~A \cite{Quillen}, \cite[IV.3.7]{Kbook}: if every
$b\backslash u$ is weakly contractible, then $u$ is a weak equivalence.

For a functor $L\colon A\to\Set$ the \emph{category of elements} $\El(L)$ has
as objects the pairs $(a,l)$ with $l\in L(a)$, and as arrows
$(a,l)\to(a',l')$ the arrows $f\colon a\to a'$ with $L(f)l=l'$. An
\emph{action component} of $A$ is a connected component of $\El(L)$ for some
$L$. If $D$ is an action component of $A$, a component of $\El(L)$, and
$K\colon D\to\Set$ is a functor, then
\begin{equation}\label{eq:composite}
 \El(K)\cong\El(\widetilde K),\qquad
 \widetilde K(a)=\coprod_{l\in L(a),\ (a,l)\in D}K(a,l) ;
\end{equation}
so every action component of $D$ is an action component of $A$. If
$\Psi\colon D\to\Set$ sends all arrows to bijections, then
$B\El(\Psi)\to BD$ is a covering space, and every covering space of $BD$
arises in this way \cite{Quillen}, \cite[IV.3.5.1 and Exercise~IV.3.1]{Kbook}.

For a functor $N\colon A\to\Ab$ let $C_\bullet(A;N)$ be the chain complex with
$C_n(A;N)=\bigoplus_{a_0\to\dots\to a_n}N(a_0)$, the sum being over the
$n$-simplices of $NA$, and with the usual alternating differential; let
$H_n(A;N)$ be its homology. We use three properties: $H_0(A;N)=\colim_AN$; a
short exact sequence of functors gives a long exact sequence of the groups
$H_n$; and for $L\colon A\to\Set$ the complex $C_\bullet(A;\Z[L])$ is the
complex of integral chains of $N\El(L)$. In particular, for
$P_x=\Z[A(x,-)]$,
\begin{equation}\label{eq:rep}
 H_0(A;P_x)=\Z,\qquad H_n(A;P_x)=0\quad(n>0),
\end{equation}
because $\El(A(x,-))=x\backslash A$ has an initial object; the isomorphism
$H_0(A;P_x)\cong\Z$ is induced by the augmentation
$\sum n_if_i\mapsto\sum n_i$.

\subsection{Definitions}

\begin{definition}\label{def:classes}
Let $A$ be a non-empty connected category.
\begin{enumerate}
\item[\textup{(a)}] $A$ is \emph{filtered} if any two objects have a common
 successor and for every pair $f,g\colon p\to q$ there is an arrow $h$ with
 $hf=hg$.
\item[\textup{(b)}] $A$ is \emph{exact} if the functor $\colim_A\colon\Ab^A\to\Ab$ is
 exact.
\item[\textup{(c)}] $A$ is \emph{total} if for every final functor $u\colon A\to B$
 the category $B$ is weakly contractible.
\item[\textup{(d)}] $A$ is \emph{quasi-filtered} if \textup{(QF1)} every pair of
 arrows $q\leftarrow p\to q'$ can be completed to a commutative square, and
 \textup{(QF2)} for every pair $f,g\colon p\to q$ there are an object $r$ and
 an element $t$ of $\Z[A(q,r)]$ of augmentation $1$ with $tf=tg$ in
 $\Z[A(p,r)]$.
\end{enumerate}
\end{definition}

Exact categories were characterized by Isbell and Mitchell
\cite{IM73,IM76}; we call an element $t$ as in (QF2) an \emph{affine
equalizer} of $f$ and $g$. (QF1) and (QF2) are the conditions I and IV of
\cite[\S 2]{IM76} for the representable functors; in \cite[p.~369]{AKR} a
category is called quasi-filtered (after Ad\'amek and Reiterman) if every
parallel pair has a quasi-filter and any two objects have a common successor,
and by \cite[\AKRone]{AKR} and \cite[Lemma~2.1]{IM76} a connected
category is quasi-filtered in that sense if and only if it satisfies (QF1)
and (QF2).
Property (c) is Grothendieck's \emph{asph\'ericit\'e totale}, for the weak
homotopy equivalences \cite[Chapter~II, ms.~p.~36ter; Chapter~XII,
D\'efinition~1]{Der}. It is not the property called totally aspherical in
the theory of test categories and in \cite[Remark~4.5(1)]{Rosicky}, which concerns the
diagonal $A\to A\times A$; see Proposition~\ref{prop:sifted}.

\subsection{Total categories}

\begin{lemma}\label{lem:acyclic}
Let $A$ be a category all of whose action components have the integral
homology of a point. Then all action components of $A$ are weakly
contractible.
\end{lemma}

\begin{proof}
Let $D$ be an action component and $\gamma$ an element of the fundamental
group of $BD$. The connected covering space of $BD$ whose fundamental group
is the cyclic group generated by $\gamma$ is $B\El(\Psi)$ for some
$\Psi\colon D\to\Set$, and $\El(\Psi)$ is an action component of $A$ by
\eqref{eq:composite}. Its first homology group is that cyclic group, so
$\gamma=1$. Thus $BD$ is simply connected with the homology of a point, and
it is weakly contractible by the Hurewicz theorem.
\end{proof}

\begin{proposition}\label{prop:total}
A non-empty connected category $A$ is total if and only if every action
component of $A$ is weakly contractible.
\end{proposition}

\begin{proof}
Suppose that all action components are weakly contractible, and let
$u\colon A\to B$ be final. For an object $b$, the category $b\backslash u$ is
the category of elements of the functor $a\mapsto B(b,u(a))$, and it is
connected; so it is an action component and is weakly contractible. By
Theorem~A, $u$ is a weak equivalence. Since $A$ is the category of elements
of the functor with value a point, $BA$ is weakly contractible, and so is
$BB$.

Conversely let $A$ be total, and let $D$ be an action component, a component
of $\El(L)$. Let $L_D\subseteq L$ be the subfunctor with
$L_D(a)=\{l:(a,l)\in D\}$. Let $A_D$ be the category obtained from $A$ by
adjoining one object $v$, with $A_D(v,a)=L_D(a)$, $A_D(a,v)=\varnothing$,
$A_D(v,v)=\{\id\}$, and with composition given by the action. The inclusion
$A\to A_D$ is final: for an object $a$ of $A$ the category $a\backslash A$
has an initial object, and $v\backslash A=D$. Hence $A_D$ is weakly
contractible, and so is $A$, because the identity of $A$ is final. A simplex
of $NA_D$ that is not in $NA$ is a chain that begins with $v$; so $NA_D$ is
the union of $NA$ and of the nerve of the category $D^{\triangleleft}$
obtained from $D$ by adjoining an initial object, glued along $ND$ by the
projection $ND\to NA$. Hence $BA_D$ is the mapping cone of $BD\to BA$, and
the reduced integral homology of $BD$ vanishes. By Lemma~\ref{lem:acyclic}
all action components are weakly contractible.
\end{proof}

\begin{remark}\label{rem:hocolim}
Equivalently, $A$ is total if and only if for every functor $X$ from $A$ to
simplicial sets the canonical map $\operatorname{hocolim}_AX\to\colim_AX$ is a
weak equivalence. Indeed, the homotopy colimit is the diagonal of the
bisimplicial set $(m,n)\mapsto N_n\El(X_m)$ \cite[Chapter~XII]{BK}, \cite[Sections~21 and~22]{JardineL8}, and the map
is induced by the maps $N\El(X_m)\to\pi_0\El(X_m)=\colim_AX_m$. If the action
components are weakly contractible these are weak equivalences, and then so is
the map of diagonals \cite[Lemma~21.1]{JardineL8}, \cite[Chapter~IV]{GJ}. Conversely a functor
$L\colon A\to\Set$ is a diagram of discrete simplicial sets, for which the map
is $N\El(L)\to\pi_0\El(L)$.
\end{remark}

\subsection{The inclusions}

\begin{proposition}\label{prop:incl}
Let $A$ be a non-empty connected category. If $A$ is filtered it is exact,
and if $A$ is exact it is total.
\end{proposition}

\begin{proof}
Filtered colimits of abelian groups are exact
\cite[Section~2.6]{Weibel}. Let $A$ be exact, and let $D$ be an action
component, a component of $\El(L)$. For $N\colon D\to\Ab$ put
$(\Sigma N)(a)=\bigoplus_{l:(a,l)\in D}N(a,l)$. The functor
$\Sigma\colon\Ab^D\to\Ab^A$ is exact and $\colim_A\Sigma N=\colim_DN$, so
$\colim_D$ is exact. Let $R_\bullet$ be the complex of functors $D\to\Ab$
with $R_n=\bigoplus_{d_0\to\dots\to d_n}\Z[D(d_n,-)]$ whose value at an object
$d$ is the complex of integral chains of the nerve of $D/d$. As $D/d$ has a
terminal object, $R_\bullet$ is a resolution of the constant functor $\Z$.
Its colimit over $D$ is the complex of integral chains of $ND$. Since
$\colim_D$ is exact, $BD$ has the integral homology of a point. By
Lemma~\ref{lem:acyclic} and Proposition~\ref{prop:total}, $A$ is total.
\end{proof}

\begin{lemma}\label{lem:succ}
Let $D$ be a connected category such that $H_1(B\El(K);\Z)=0$ for every
functor $K\colon D\to\Set$ with $\El(K)$ connected. Then any two objects of
$D$ have a common successor.
\end{lemma}

\begin{proof}
Let $x,x'$ be objects, and let $U$ and $U'$ be the full subcategories of the
objects that receive an arrow from $x$, respectively from $x'$. They are
cosieves. Suppose that they are disjoint and put $V=U\cup U'$. Let
$K(d)=\{1,2\}$ for $d\notin V$ and $K(d)=\{1\}$ for $d\in V$, with the
evident maps. Then $N\El(K)$ is the union of two copies of $ND$ that meet in
$NV$, and $\El(K)$ is connected. In the Mayer--Vietoris sequence
\[
 H_1(B\El(K);\Z)\longrightarrow H_0(BV;\Z)\longrightarrow
 H_0(BD;\Z)\oplus H_0(BD;\Z)
\]
the middle group has rank at least $2$, since no arrow joins $U$ and $U'$,
and the image of the second map has rank $1$. So the first group is not
zero, a contradiction.
\end{proof}

\begin{proposition}\label{prop:qf}
Let $A$ be a non-empty connected category such that $H_1(BD;\Z)=0$ for every
action component $D$ of $A$. Then $A$ is quasi-filtered. In particular every
total category is quasi-filtered.
\end{proposition}

\begin{proof}
By \eqref{eq:composite}, Lemma~\ref{lem:succ} applies to every action
component of $A$.

(QF1). Let $f\colon p\to q$ and $f'\colon p\to q'$. The category
$p\backslash A$ is an action component. A common successor of its objects
$(q,f)$ and $(q',f')$ is a commutative square $uf=u'f'$.

(QF2). Let $f,g\colon p\to q$. Let $H$ be the quotient of the functor
$A(p,-)$ by the equivalence relation generated by $kf\sim kg$ for all arrows
$k$ with source $q$. Then $\El(H)$ is connected, because every object
receives an arrow from $(p,[\id_p])$. There is an exact sequence
\[
 P_q\xrightarrow{\ \partial\ }P_p\longrightarrow\Z[H]\longrightarrow 0,
 \qquad\partial(k)=kf-kg .
\]
Let $K=\ker\partial$ and $J=\operatorname{im}\partial$. The sequences
$0\to J\to P_p\to\Z[H]\to0$ and $0\to K\to P_q\to J\to 0$ and \eqref{eq:rep}
give
\[
 H_1(A;\Z[H])=\ker\bigl(\colim J\to\colim P_p\bigr),\qquad
 \colim J=\operatorname{coker}\bigl(\colim K\to\colim P_q=\Z\bigr).
\]
The elements of $J$ have augmentation zero, so $\colim J\to\colim P_p=\Z$ is
zero. Hence $H_1(B\El(H);\Z)=H_1(A;\Z[H])\cong\Z/I$, where $I$ is the subgroup
of $\Z$ generated by the augmentations of the elements of the groups
$K(x)=\{t\in\Z[A(q,x)]:tf=tg\}$. By hypothesis $I=\Z$: there are
$t_j\in K(x_j)$ and integers $n_j$ with $\sum_jn_j\,\varepsilon(t_j)=1$. By
Lemma~\ref{lem:succ}, applied to $A$, there are arrows $u_j\colon x_j\to r$
into one object $r$, and $t=\sum_jn_ju_jt_j$ has augmentation $1$ and
satisfies $tf=tg$.

The last assertion follows from Proposition~\ref{prop:total}.
\end{proof}

\subsection{Sifted categories}

A category $A$ is \emph{sifted} if it is non-empty and the diagonal
$\Delta\colon A\to A\times A$ is final, that is, if for all objects $a,b$ the
category $(a,b)\backslash\Delta$ of the pairs of arrows $a\to c\leftarrow b$ is
non-empty and connected. By \cite[Theorem~1.6]{AR} this holds if and only if
colimits over $A$ of diagrams of sets commute with finite products, which is
the definition of sifted categories in \cite[Definition~1.1]{AR}.

\begin{proposition}\label{prop:sifted}
A non-empty connected category is filtered if and only if it is total and
sifted.
\end{proposition}

\begin{proof}
Let $A$ be total and sifted. For objects $a,b$ the category
$(a,b)\backslash\Delta$ is non-empty, so $a$ and $b$ have a common successor.
Let $f,g\colon p\to q$. The category $(p,p)\backslash\Delta$ is the category
of elements of $A(p,-)\times A(p,-)$, and it is connected; so it is an action
component of $A$. By Proposition~\ref{prop:total} and Lemma~\ref{lem:succ}
its objects $(q,f,g)$ and $(q,f,f)$ have a common successor $(r,w,w')$:
there are $h,h'\colon q\to r$ with $hf=w=h'f$ and $hg=w'=h'f$. Hence $hf=hg$,
and $A$ is filtered.

Conversely a filtered category is total by Proposition~\ref{prop:incl}, and
it is sifted: two objects $(c,u,v)$ and $(c',u',v')$ of
$(a,b)\backslash\Delta$ are joined through an object $(d,\cdot,\cdot)$, where
$c\to d\leftarrow c'$ are chosen so as to equalize the two resulting arrows
$a\to d$ and the two resulting arrows $b\to d$.
\end{proof}
% ---- end sec2_classes.tex ----
% ---- begin sec3_reduction.tex (inlined for the deposit) ----
\section{A reduction for categories of embeddings}\label{sec:reduction}

Both examples of this paper are categories of finite substructures of a
homogeneous structure. This section isolates what is used of such a category.

\subsection{The setting}

Let $M$ be a non-empty set and $G$ a group of permutations of $M$. For
$X\subseteq M$ let $G_X$ be the pointwise stabilizer of $X$; so
$G_\varnothing=G$. A subgroup of $G$ is \emph{open} if it contains $G_X$ for
some finite $X\subseteq M$. Let $\K$ be a set of finite subsets of $M$. We
consider the following conditions.
\begin{enumerate}
\item[\textup{(K1)}] $\varnothing\in\K$; $gB\in\K$ for $B\in\K$ and $g\in G$; every
 finite subset of $M$ is contained in a member of $\K$.
\item[\textup{(K2)}] $\Fix_M(G_B)=B$ for every $B\in\K$.
\item[\textup{(K3)}] Every open subgroup of $G$ is of the form $G_F$ with $F\in\K$.
\end{enumerate}
Let $\mathcal A=\mathcal A(M,G,\K)$ be the category whose objects are the
non-empty members of $\K$ and whose arrows $B\to C$ are the maps
$g|_B\colon B\to C$ with $g\in G$ and $gB\subseteq C$, composed as maps.
Inclusions between members of $\K$ are arrows. By (K1) any two objects are
contained in a third, so $\mathcal A$ is non-empty and connected.

For an object $P$ of a category $\mathcal C$, the objects of
$P\backslash\mathcal C$ are the pairs $(X,x\colon P\to X)$, and an arrow
$(X,x)\to(Y,y)$ is an arrow $f\colon X\to Y$ with $fx=y$. A \emph{cosieve} of a
category is a full subcategory that contains the target of every arrow whose
source it contains.
\begin{enumerate}
\item[\textup{(SC)}] For every object $P$ of $\mathcal C$, every non-empty cosieve
 of $P\backslash\mathcal C$ is weakly contractible.
\item[\textup{(K4)}] The category $\mathcal A(M,G,\K)$ satisfies \textup{(SC)}.
\end{enumerate}
Property (SC) is invariant under equivalence of categories: an equivalence
$\Phi\colon\mathcal C\to\mathcal C'$ induces equivalences
$P\backslash\mathcal C\to\Phi P\backslash\mathcal C'$, and the objects
isomorphic to images of objects of a cosieve $\mathcal U$ form a cosieve
equivalent to $\mathcal U$.

\begin{lemma}\label{lem:squares}
In a category with \textup{(SC)} every pair of arrows $f\colon a\to b$,
$f'\colon a\to b'$ can be completed to a commutative square $uf=u'f'$.
\end{lemma}

\begin{proof}
In $a\backslash\mathcal C$ let $\mathcal U$ be the cosieve of the objects
$(c,h)$ such that $h$ factors through $f$, and $\mathcal U'$ the one defined by
$f'$. Their union is a non-empty cosieve, hence connected. A path in it from
$(b,f)$ to $(b',f')$ contains an arrow between an object of $\mathcal U$ and an
object of $\mathcal U'$, and its target lies in $\mathcal U\cap\mathcal U'$.
\end{proof}

\subsection{The components of the categories of elements}

\begin{proposition}\label{prop:reduction}
Assume \textup{(K1)--(K4)}, and let $L\colon\mathcal A\to\Set$ be a functor. For
every connected component $D$ of $\El(L)$, either $D$ is weakly contractible,
or the projection $D\to\mathcal A$ is a weak equivalence.
\end{proposition}

\begin{proof}
A functor that induces a weak equivalence of classifying spaces induces a
bijection between the sets of connected components and weak equivalences
between corresponding components. The two functors constructed below commute
with the projections to $\mathcal A$.

\emph{Step 1: injective transition maps.} For $l,l'\in L(a)$ write $l\sim l'$
if $L(f)l=L(f)l'$ for some arrow $f\colon a\to b$. This is an equivalence
relation: if $f$ and $f'$ witness $l\sim l'$ and $l'\sim l''$, and $uf=u'f'$ as
in Lemma~\ref{lem:squares}, then $L(uf)$ takes the same value on $l$, $l'$ and
$l''$. It is preserved by $L(v)$ for every arrow $v\colon a\to a'$: if
$L(f)l=L(f)l'$ and $uf=u'v$, then $L(u')$ identifies $L(v)l$ and $L(v)l'$.
So the quotients form a functor $L^{\mathrm{inj}}$, with a natural surjection
$L\to L^{\mathrm{inj}}$. Its transition maps are injective: an arrow that
witnesses $L(v)l\sim L(v)l'$, composed with $v$, witnesses $l\sim l'$.

We claim that $q\colon\El(L)\to\El(L^{\mathrm{inj}})$ is a weak equivalence. Fix
an object $(a,[l])$ of the target and let $\mathcal D=(a,[l])\backslash q$; its
objects are the triples $(x,k,f)$ with $f\colon a\to x$, $k\in L(x)$ and
$k\sim L(f)l$. Let $J\colon a\backslash\mathcal A\to\mathcal D$ send
$f\colon a\to x$ to $(x,L(f)l,f)$. For an object $d=(x,k,f)$ of $\mathcal D$,
the category $d\backslash J$ is the full subcategory of $x\backslash\mathcal A$
on the arrows $g\colon x\to y$ with $L(g)k=L(gf)l$. It is a cosieve, and it is
non-empty because $k\sim L(f)l$; by (K4) it is weakly contractible. By
Quillen's Theorem~A (Section~\ref{sec:classes}), $J$ is a weak equivalence. As
$a\backslash\mathcal A$ has an initial object, $\mathcal D$ is weakly
contractible, and Theorem~A applies to $q$.

\emph{Step 2: saturation.} By Step 1 we may assume that all maps $L(f)$ are
injective. The non-empty members of $\K$ form a directed set under inclusion,
by (K1). Let
\[
 S=\colim_{B}L(B)
\]
be the colimit over the inclusions. The canonical maps $i_B\colon L(B)\to S$
are injective. The group $G$ acts on $S$ by
$g\cdot i_B(l)=i_{gB}\bigl(L(g|_B)l\bigr)$, where $g|_B\colon B\to gB$; this is
well defined because $g|_{B'}$ extends $g|_B$ when $B\subseteq B'$. The element
$i_B(l)$ is fixed by $G_B$. So $S$ is a $G$-set all of whose point stabilizers
are open.

For any such $G$-set $S$ define a functor $T_S\colon\mathcal A\to\Set$ by
$T_S(B)=S^{G_B}$ and $T_S(g|_B)(s)=g\cdot s$ for an arrow $g|_B\colon B\to C$.
This does not depend on $g$, since two choices differ by an element of $G_B$;
and $g\cdot s$ is fixed by $G_C$ because $g^{-1}G_Cg\le G_B$. In our case the
maps $i_B$ form a natural monomorphism $i\colon L\to T_S$. For an object
$(a,s)$ of $\El(T_S)$, the category $(a,s)\backslash\El(i)$ is the full
subcategory of $a\backslash\mathcal A$ on the arrows $f\colon a\to b$ with
$T_S(f)s\in i_b(L(b))$. It is a cosieve. It is non-empty: if $s=i_c(l)$,
choose $b\in\K$ containing $a\cup c$; for the inclusion $f\colon a\to b$ we
have $T_S(f)s=s\in i_b(L(b))$. By (K4) and Theorem~A,
$\El(i)\colon\El(L)\to\El(T_S)$ is a weak equivalence.

\emph{Step 3: orbits.} The functor $S\mapsto T_S$ takes disjoint unions of
$G$-sets to disjoint unions of functors, so $\El(T_S)$ is the disjoint union of
the categories $\El(T_O)$, $O$ running through the orbits of $S$. By (K3)
every orbit is isomorphic to $G/G_F$ for some $F\in\K$.

If $F\ne\varnothing$, then $T_{G/G_F}\cong\mathcal A(F,-)$. Indeed, a coset
$gG_F$ is fixed by $G_B$ if and only if $G_B\le gG_Fg^{-1}=G_{gF}$, that is, by
(K2), if and only if $gF\subseteq B$. The map $gG_F\mapsto g|_F$ is then a
bijection from $(G/G_F)^{G_B}$ onto $\mathcal A(F,B)$, natural in $B$. The
category of elements of $\mathcal A(F,-)$ is $F\backslash\mathcal A$, which has
an initial object.

If $F=\varnothing$, the orbit is a point, $T_O$ is the constant functor with
value a point, and $\El(T_O)\to\mathcal A$ is an isomorphism.

So every component of $\El(T_S)$ is contractible or is mapped isomorphically
to $\mathcal A$ by the projection, and the statement follows from Steps 1
and 2.
\end{proof}

\subsection{The bridge set}

The following construction is used twice, to verify \textup{(K3)}. It involves
only finite linearly ordered sets.

\begin{lemma}\label{lem:bridgeset}
Let $D$ and $E$ be finite linearly ordered sets, $B,C\subseteq D$,
$F=B\cap C$, and let $d_0,d_1\colon D\to E$ be increasing injections that agree
on $F$. For every sufficiently large integer $N\ge 2$, with $s=2N$, there are a
finite linearly ordered set $Z$, an increasing injection $e\colon E\to Z$ and
increasing injections $a_0,\dots,a_s\colon D\to Z$ with the following
properties. Put $Z_0=e(E)$ and $Z_i=a_i(D)$ for $0<i<s$; these are the
\emph{pieces}, and their indices are read modulo $s$.
\begin{enumerate}
\item[\textup{(i)}] $Z=\bigcup_iZ_i$; $a_0=ed_0$ and $a_s=ed_1$;
 $a_{2j}|_B=a_{2j+1}|_B$ and $a_{2j+1}|_C=a_{2j+2}|_C$ for $0\le j<N$.
\item[\textup{(ii)}] For $x\in F$ the point $a_i(x)$ does not depend on $i$; these
 \emph{$F$-points} lie in every piece. Every other point of $Z$ lies in at
 most three pieces. If it lies in two pieces, these are consecutive, or they
 are $Z_1$ and $Z_{s-1}$, and then the point lies in $Z_0$ as well.
\item[\textup{(iii)}] For $0\le i<s$ we have $Z_i\cap Z_{i+1}=a_i(B)=a_{i+1}(B)$ if
 $i$ is even and $Z_i\cap Z_{i+1}=a_i(C)=a_{i+1}(C)$ if $i$ is odd.
\item[\textup{(iv)}] For every index $r$ there are a finite set $L$, a
 surjection $q\colon L\to Z$ and subsets $T_0,\dots,T_{s-1}$ covering $L$ such
 that $q$ maps $T_k$ bijectively onto $Z_{r+k}$, with these properties: every
 point of $L$ lies in $T_k$ for a set of consecutive values of $k$, with at
 most three elements unless the point lies over an $F$-point;
 $q(T_{k-1}\cap T_k)=Z_{r+k-1}\cap Z_{r+k}$ for $0<k<s$; every $F$-point has
 one preimage under $q$; and $q$ is injective except that at most
 $\max(|B|,|C|)$ points of $Z$ have two preimages.
\end{enumerate}
\end{lemma}

\begin{proof}
Let $\widehat Z$ be the disjoint union of $E$ and of copies $D_1,\dots,D_{s-1}$
of $D$; write $x_i$ for the copy of $x\in D$ in $D_i$. We call $E$ the
\emph{piece} $0$ and $D_i$ the piece $i$. Consider the following \emph{links}
between points of consecutive pieces: $d_0(x)$ and $x_1$, for $x\in B$; $x_i$
and $x_{i+1}$, for $0<i<s-1$ and $x\in B$ if $i$ is even, $x\in C$ if $i$ is
odd; $x_{s-1}$ and $d_1(x)$, for $x\in C$. Let $Z$ be the quotient of
$\widehat Z$ by the equivalence relation generated by the links, $e$ the
restriction of the quotient map to $E$, $a_i(x)$ the class of $x_i$ for
$0<i<s$, and $a_0=ed_0$, $a_s=ed_1$. Then (i) holds.

A point has at most one link to the next piece and at most one to the
preceding piece, because $d_0$ and $d_1$ are injective. For $x\in F$ the
points $d_0(x)=d_1(x),x_1,\dots,x_{s-1}$ form a class, with one point in each
piece. For $x\in B\setminus C$ the links at copies of $x$ join $x_i$ to
$x_{i+1}$ for $i$ even, $0<i<s-1$, and $d_0(x)$ to $x_1$; for
$x\in C\setminus B$ they join $x_i$ to $x_{i+1}$ for $i$ odd, $0<i<s-1$, and
$x_{s-1}$ to $d_1(x)$; for $x\notin B\cup C$ there are none. A point of $E$ that
is not of the form $d_0(x)$, $x\in F$, has at most two links, to some $x_1$
with $x\in B\setminus C$ and to some $y_{s-1}$ with $y\in C\setminus B$. Hence a
class that does not come from $F$ is a single point, or a pair
$\{x_i,x_{i+1}\}$, or is contained in a set $\{y_{s-1},d,x_1\}$ with $d\in E$.
So every class has at most one point in each piece, which shows that $e$ and
the $a_i$ are injective, and (ii) and (iii) follow by inspection; here
$N\ge 2$ is used.

For (iv), let $L$ be the quotient of $\widehat Z$ by the relation generated
by the links other than those between the pieces $r-1$ and $r$, let $T_k$ be
the image of the piece $r+k$, and $q$ the quotient map. A class of $Z$ is the
union of at most two classes of $L$, and of two only if it contains a link
between the pieces $r-1$ and $r$ and does not come from $F$; there are at
most $\max(|B|,|C|)$ such links. The remaining assertions follow from the
description of the classes.

\emph{The order.} Choose an increasing injection $\varepsilon\colon E\to\R$
and put $p_t=(1-t)\,\varepsilon d_0+t\,\varepsilon d_1\colon D\to\R$ for
$0\le t\le 1$. Each $p_t$ is strictly increasing, and $p_t(x)$ does not
depend on $t$ for $x\in F$. Let $\Delta>0$ be a lower bound for
$p_t(x')-p_t(x)$ over all $t$ and all $x<x'$ in $D$, and let
$M_0=\max_x|\varepsilon d_1(x)-\varepsilon d_0(x)|$; assume $2M_0<N\Delta$.
Define $b_i\colon D\to\R$ for $0\le i\le s$ by
\[
 b_{2j}=p_{j/N},\qquad
 b_{2j+1}(x)=\begin{cases}p_{j/N}(x),&x\in B,\\ p_{(j+1)/N}(x),&x\notin B.\end{cases}
\]
Each value of $b_{2j+1}$ differs from the corresponding value of $p_{(j+1)/N}$
by at most $M_0/N$, so $b_{2j+1}(x')-b_{2j+1}(x)\ge\Delta-2M_0/N>0$ for
$x<x'$: all $b_i$ are strictly increasing. Moreover $b_0=\varepsilon d_0$,
$b_s=\varepsilon d_1$, $b_{2j}$ and $b_{2j+1}$ agree on $B$, and $b_{2j+1}$
and $b_{2j+2}$ agree on $C$ (on $C\setminus B$ by definition, on $F$ because
$p_t$ is constant there). So the function on $\widehat Z$ that is
$\varepsilon$ on $E$ and $b_i$ on $D_i$ takes the same value at the two ends
of every link, and induces a function $\beta\colon Z\to\R$ that is injective
and increasing on every piece. Any linear order on $Z$ in which $z<z'$
whenever $\beta(z)<\beta(z')$ makes $e$ and all $a_i$ increasing.
\end{proof}
% ---- end sec3_reduction.tex ----
% ---- begin sec4_cdelta.tex (inlined for the deposit) ----
\section{A total category whose colimits are not exact}\label{sec:cdelta}

\subsection{The category $\Cd$}

A \emph{hypergraph} is a finite set $X$ (of vertices) with a set $E_X$ of
three-element subsets of $X$ (the hyperedges). A subset $U\subseteq X$
carries the induced hypergraph, with hyperedges
$E_X(U)=\{h\in E_X: h\subseteq U\}$. Put
\[
 \delta_X(U)=2\,|U|-3\,|E_X(U)| .
\]
We write $\delta(U)$ when $X$ is understood; $\delta(U)$ depends only on the
induced hypergraph on $U$. The hypergraph $X$ is \emph{positive} if
$\delta(U)>0$ for every non-empty $U\subseteq X$. For $A\subseteq B\subseteq X$
we say that $A$ is \emph{closed in} $B$, and write $A\cle B$, if
\[
 \delta(U)>\delta(A)\qquad\text{for every } U \text{ with } A\subsetneq U\subseteq B .
\]
Thus $X$ is positive if and only if $\varnothing\cle X$. An \emph{ordered}
hypergraph carries in addition a linear order on its vertices. An
\emph{embedding} $f\colon X\to Y$ of ordered hypergraphs is an increasing
injection such that a three-element subset $h$ of $X$ is a hyperedge of $X$
exactly when $f(h)$ is a hyperedge of $Y$; it is a \emph{closed embedding} if
$f(X)\cle Y$.

\begin{definition}\label{def:cdelta}
The objects of $\Cd$ are the positive ordered hypergraphs with vertex set
$[n]=\{1<\dots<n\}$, $n\ge 1$. Its arrows are the closed embeddings, composed as
maps.
\end{definition}

Composites of closed embeddings are closed by Lemma~\ref{lem:calc}(c) below, so
$\Cd$ is a category; it is small, and every non-empty positive ordered
hypergraph is isomorphic to exactly one object. We write $V$ for the object
with one vertex and $W$ for the object $([3],\{\{1,2,3\}\})$. The latter is
positive: its subsets with one, two and three elements have
$\delta=2,4,3$. The same numbers show that the three maps
$v_1,v_2,v_3\colon V\to W$, $v_i(1)=i$, are closed embeddings.

\begin{lemma}[closure calculus]\label{lem:calc}
Let $X$ be a hypergraph.
\begin{enumerate}
\item[\textup{(a)}] $\delta(U)+\delta(U')\ge\delta(U\cup U')+\delta(U\cap U')$ for all
 $U,U'\subseteq X$.
\item[\textup{(b)}] If $A\cle B$ and $Y\subseteq B$, then $A\cap Y\cle Y$.
\item[\textup{(c)}] If $A\cle B$ and $B\cle C$, then $A\cle C$.
\item[\textup{(d)}] Let $X=B\cup C$ and $A=B\cap C$, and suppose that
 $E_X=E_X(B)\cup E_X(C)$ \textup{(}we then call $X$ the \emph{free amalgam} of $B$
 and $C$ over $A$\textup{)}. Then for every $U\subseteq X$
 \begin{equation}\label{eq:amalgam}
  \delta(U)=\delta(U\cap B)+\delta(U\cap C)-\delta(U\cap A).
 \end{equation}
 If moreover $B$ and $C$ are positive, $A\cle B$ and $A\cle C$, then $X$ is
 positive, $B\cle X$ and $C\cle X$. Here $A$ may be empty.
\end{enumerate}
\end{lemma}

\begin{proof}
(a) The number of vertices is modular. For the hyperedges,
$E(U)\cup E(U')\subseteq E(U\cup U')$ and $E(U)\cap E(U')=E(U\cap U')$, so
$|E(U)|+|E(U')|\le|E(U\cup U')|+|E(U\cap U')|$.

(b) Let $A\cap Y\subsetneq U\subseteq Y$. Then $U\cap A=A\cap Y$ and
$U\cup A\supsetneq A$, so by (a)
$\delta(U)-\delta(A\cap Y)\ge\delta(U\cup A)-\delta(A)>0$.

(c) Let $A\subsetneq U\subseteq C$. If $U\subseteq B$ there is nothing to prove.
Otherwise $U\cap B\subsetneq U$ and $U\cap B\cle U$ by (b), so
$\delta(U)>\delta(U\cap B)\ge\delta(A)$, the last inequality because
$A\subseteq U\cap B\subseteq B$ and $A\cle B$.

(d) A hyperedge of $X$ lies in $B$ or in $C$, and those lying in both are the
hyperedges contained in $A$; this gives \eqref{eq:amalgam}. Let $U$ be
non-empty. If $U\subseteq B$ then $\delta(U)>0$. Otherwise
$U\cap A\subsetneq U\cap C$, and $U\cap A\cle U\cap C$ by (b), so
$\delta(U\cap C)-\delta(U\cap A)>0$, while $\delta(U\cap B)\ge 0$; hence
$\delta(U)>0$. If $B\subsetneq U\subseteq X$, then $U\cap B=B$, $U\cap A=A$ and
$U\cap C\supsetneq A$, so
$\delta(U)-\delta(B)=\delta(U\cap C)-\delta(A)>0$. Thus $B\cle X$, and
$C\cle X$ by symmetry.
\end{proof}

If, in (d), $B$ and $C$ are ordered and their orders agree on $A$, the free
amalgam carries linear orders extending both. We call \emph{the ordered free
amalgam with $B$ first} the one in which, for $b\in B\setminus A$ and
$c\in C\setminus A$, we have $b<c$ exactly when every element of $A$ below $b$
is below $c$.

\subsection{The generic structure and its automorphism group}

For an ordered hypergraph $M$ with countably many vertices, all of whose finite
subsets are positive, and a finite $B\subseteq M$, we write $B\cle M$ if
$B\cle U$ for every finite $U$ with $B\subseteq U\subseteq M$. Parts (b) and
(c) of Lemma~\ref{lem:calc} remain true with $M$ in the place of the largest
set.

\begin{proposition}\label{prop:generic}
There is a countable ordered hypergraph $M$, all of whose finite subsets are
positive, with the following properties.
\begin{enumerate}
\item[\textup{(M1)}] Every finite subset of $M$ is contained in a finite subset
 $B$ with $B\cle M$.
\item[\textup{(M2)}] If $B\cle M$ is finite and $f\colon B\to C$ is a closed
 embedding into a finite positive ordered hypergraph, there is an embedding
 $h\colon C\to M$ with $h(C)\cle M$ and $hf=\id_B$.
\item[\textup{(M3)}] Every isomorphism between two finite closed subsets of $M$
 extends to an automorphism of $M$.
\end{enumerate}
\end{proposition}

\begin{proof}
A \emph{task} is a triple $(k,B,f)$ where $k\ge 0$ and $f\colon B\to C$ is a
closed embedding of a subset $B$ of the $k$-th stage, closed in that stage,
into a finite positive ordered hypergraph. Up to isomorphism of $C$ under $B$
there are countably many tasks at each stage, so the tasks of all stages can
be enumerated in such a way that every task is met at a later stage. Starting
with $M_0=V$, let $M_{l+1}$ be an ordered free amalgam of $M_l$ and $C$ over
$B$, where $(k,B,f)$ is the task met at stage $l\ge k$. Here $B\cle M_k\cle M_l$
by induction and Lemma~\ref{lem:calc}(c), so by Lemma~\ref{lem:calc}(d) the
amalgam is positive, $M_l\cle M_{l+1}$, and the image of $C$ is closed in
$M_{l+1}$. Let $M=\bigcup_l M_l$. By transitivity $M_l\cle M_{l'}$ for
$l\le l'$, hence $M_l\cle M$; this gives (M1). For (M2), choose $k$ with
$B\subseteq M_k$; then $B\cle M_k$, the task
$(k,B,f)$ is met at some stage $l$, and the resulting copy of $C$ is closed in
$M_{l+1}\cle M$, hence in $M$. For (M3) one extends a given isomorphism
$\varphi\colon B\to B'$ of finite closed subsets step by step: for $x\in M$
choose by (M1) a finite closed $C\supseteq B\cup\{x\}$; by (M2), applied to
$B'$ and to $\varphi^{-1}\colon B'\to B\subseteq C$, the map $\varphi$ extends to
an isomorphism of $C$ onto a closed subset of $M$. Alternating between
$\varphi$ and $\varphi^{-1}$ along an enumeration of $M$ gives an automorphism.
\end{proof}

The existence of such limits for classes with a notion of strong
substructure is standard; see
\cite[Theorem~2.6]{EHN}. We fix $M$ as in Proposition~\ref{prop:generic} and put
\[
 G=\Aut(M),\qquad \K=\{B\subseteq M:\ B \text{ finite},\ B\cle M\},
\]
where automorphisms preserve the order and the hyperedges. Positivity gives
$\varnothing\in\K$. We use the notation of Section~\ref{sec:reduction}:
$G_B$ is the pointwise stabilizer of $B$, and $\mathcal A=\mathcal A(M,G,\K)$.

\begin{lemma}\label{lem:AMequiv}
The triple $(M,G,\K)$ satisfies \textup{(K1)}. For $B,C\in\K$ the arrows
$B\to C$ of $\mathcal A$ are exactly the closed embeddings $B\to C$, and
$\mathcal A$ is equivalent to $\Cd$.
\end{lemma}

\begin{proof}
(K1) follows from (M1) and from the invariance of $\cle$ under automorphisms.
If $g\in G$ and $gB\subseteq C$, then $gB\cle M$, so $gB\cle C$ by
Lemma~\ref{lem:calc}(b), and $g|_B$ is a closed embedding. Conversely, if
$f\colon B\to C$ is a closed embedding, then $f(B)\cle C\cle M$, so
$f(B)\in\K$ by Lemma~\ref{lem:calc}(c), and $f\colon B\to f(B)$ extends to an
element of $G$ by (M3). Finally every object of $\Cd$ is isomorphic to a member
of $\K$ by (M2) with $B=\varnothing$, so the functor that sends a non-empty
$B\in\K$ to the object of $\Cd$ isomorphic to it is an equivalence.
\end{proof}

\begin{lemma}\label{lem:K2delta}
$(M,G,\K)$ satisfies \textup{(K2)}: $\Fix_M(G_B)=B$ for every $B\in\K$.
\end{lemma}

\begin{proof}
Let $x\in M\setminus B$ and choose $D\in\K$ containing $B\cup\{x\}$. Let $D'$ be a
second copy of $D$, with $D\cap D'=B$, and form an ordered free amalgam of $D$
and $D'$ over $B$; note that $B\cle D$ by Lemma~\ref{lem:calc}(b). By
Lemma~\ref{lem:calc}(d) the inclusion of $D$ into it is a closed embedding, so
by (M2) the amalgam embeds into $M$ over $D$, with closed image; the image of
$D'$ is then closed in $M$. The isomorphism of $D$ onto this image that is the
identity on $B$ extends by (M3) to some $g\in G_B$, and $gx\ne x$.
\end{proof}

\begin{lemma}\label{lem:SCdelta}
The category $\mathcal A$ satisfies \textup{(SC)}, and $B\mathcal A$ is weakly
contractible.
\end{lemma}

\begin{proof}
Both properties are invariant under equivalence, so we may replace $\mathcal A$
by the following category $\mathcal C$. Fix an infinite set $\Omega$ and
injections $\lambda,\rho\colon\Omega\to\Omega$ with disjoint images. The objects
of $\mathcal C$ are the non-empty positive ordered hypergraphs whose vertex set
is a finite subset of $\Omega$, and its arrows are the closed embeddings; by
Lemma~\ref{lem:AMequiv} it is equivalent to $\mathcal A$.

Let $P$ be an object of $\mathcal C$, let $\mathcal U$ be a non-empty full
cosieve of $P\backslash\mathcal C$, and fix an object $(Z,z\colon P\to Z)$ of
$\mathcal U$. For an object $(X,x\colon P\to X)$ of $P\backslash\mathcal C$ let
$T(X)$ be the hypergraph on $\lambda(X)\cup\rho(Z\setminus z(P))$ that is the free
amalgam of the copy $\lambda(X)$ of $X$ and of the copy
$\lambda(x(P))\cup\rho(Z\setminus z(P))$ of $Z$ over $\lambda(x(P))$, ordered
with $\lambda(X)$ first. The base is closed in both copies because $x$ and
$z$ are closed embeddings, so by Lemma~\ref{lem:calc}(d) $T(X)$ is positive and
the maps
\[
 \eta_X\colon X\to T(X),\ u\mapsto\lambda(u),\qquad
 \gamma_X\colon Z\to T(X),\ z(p)\mapsto\lambda(x(p)),\ w\mapsto\rho(w)\ \
 (w\notin z(P)),
\]
are closed embeddings; they are arrows under $P$. For an arrow
$f\colon X\to X'$ of $P\backslash\mathcal C$ let
$T(f)=\lambda f\lambda^{-1}$ on $\lambda(X)$ and the identity on
$\rho(Z\setminus z(P))$. It is increasing, because $f$ is increasing and
commutes with the maps from $P$, so that it preserves the gap of $P$ in which a
vertex lies. It is an embedding, since a hyperedge of $T(X')$ lies in one of
the two copies and $f$ is an embedding. Its image is closed, because $T(X')$
is the free amalgam of $\lambda(X')$ and of the image of $T(X)$ over
$\lambda(f(X))$, and $f(X)\cle X'$. So $T$ is an endofunctor of
$P\backslash\mathcal C$, and $\eta\colon\id\to T$ and
$\gamma\colon\mathrm{const}_{(Z,z)}\to T$ are natural transformations. Since
$\mathcal U$ is a cosieve containing $(Z,z)$, the arrows $\gamma_X$ show that
$T$ maps $P\backslash\mathcal C$ into $\mathcal U$. The restrictions of $\eta$
and $\gamma$ to $\mathcal U$ connect the identity of $\mathcal U$ to a constant
functor, so $B\mathcal U$ is contractible \cite[IV.3.2]{Kbook}.

The same construction with the empty set in the place of $P$ (so that
$T(X)=\lambda(X)\cup\rho(Z)$ is a disjoint union, with $\lambda(X)$ first)
gives, for a fixed object $Z$, an endofunctor $T$ of $\mathcal C$ with natural
transformations $\id\to T\leftarrow\mathrm{const}_Z$. Hence $B\mathcal C$ is
contractible.
\end{proof}

\subsection{Open subgroups}

\begin{lemma}[the bridge]\label{lem:bridgedelta}
Let $D$ and $E$ be positive ordered hypergraphs, let $B\cle D$ and $C\cle D$,
put $F=B\cap C$, and let $d_0,d_1\colon D\to E$ be closed embeddings that agree
on $F$. Then there are a positive ordered hypergraph $Z$, a closed embedding
$e\colon E\to Z$ and closed embeddings $a_0,\dots,a_{2N}\colon D\to Z$ with
\[
 a_0=ed_0,\quad a_{2N}=ed_1,\quad a_{2j}|_B=a_{2j+1}|_B,\quad
 a_{2j+1}|_C=a_{2j+2}|_C\qquad(0\le j<N).
\]
\end{lemma}

\begin{proof}
Put $m=\max(|B|,|C|)$ and take $N$ so large that Lemma~\ref{lem:bridgeset}
applies and
\begin{equation}\label{eq:Nbound}
 s:=2N>3\,\bigl(\delta(E)+2m\bigr),\qquad N\ge 2 .
\end{equation}
Let $Z$, $e$, $a_i$ and the pieces $Z_0,\dots,Z_{s-1}$ be as in
Lemma~\ref{lem:bridgeset}; the indices of pieces are read modulo $s$. Every
piece is a copy of $E$ or of $D$; let the hyperedges of $Z$ be the images of
the hyperedges of the pieces. It remains to prove that every piece is an
induced subhypergraph of $Z$, that $Z$ is positive, and that every piece is
closed in $Z$.

\emph{Pieces are induced.} For two consecutive pieces the intersection
$Z_i\cap Z_{i+1}$ is the common image of $B$ or of $C$ (Lemma~\ref{lem:bridgeset}(iii)), and the two pieces induce
the same hypergraph on it, because $d_0$ and $d_1$ are embeddings. Call this
fact $(\ast)$. Let $h$ be a hyperedge of the piece $Z_j$ whose three vertices lie in
$Z_i$, $i\ne j$. If all three are $F$-points, $h$ is the image of a hyperedge
of $D$ contained in $F$, hence a hyperedge of every piece. Otherwise one of
them lies in at most three cyclically consecutive pieces, among them $Z_i$ and
$Z_j$. If $i$ and $j$ are consecutive, $h$ is a hyperedge of $Z_i$ by $(\ast)$.
If not, then $\{i,j\}=\{1,s-1\}$ by Lemma~\ref{lem:bridgeset}(ii), and every
point of $Z_1\cap Z_{s-1}$ lies in $Z_0$; applying $(\ast)$ to $(Z_j,Z_0)$ and
then to $(Z_0,Z_i)$ shows again that $h$ is a hyperedge of $Z_i$.

\emph{The cut.} Fix an index $r$. Let $L$ be obtained from the disjoint union
of the pieces by making all the identifications of Lemma~\ref{lem:bridgeset}
except those between the pieces $r-1$ and $r$, with hyperedges the images of
the hyperedges of the pieces, and let $q\colon L\to Z$ be the quotient map.
Write $T_0,\dots,T_{s-1}$ for the images in $L$ of the pieces
$r,r+1,\dots,r-1$, in this order. By Lemma~\ref{lem:bridgeset}(iv) a point of
$L$ lies in $T_k$ for a set of consecutive values of $k$. Hence
$T_k\cap(T_0\cup\dots\cup T_{k-1})=T_k\cap T_{k-1}=:A_k$. Applying $(\ast)$ along
the consecutive pieces between two pieces that contain the vertices of a
hyperedge shows, as in the first step, that every $T_k$ is an induced
subhypergraph of $L$. So $L_k=T_0\cup\dots\cup T_k$ is the free amalgam of
$L_{k-1}$ and $T_k$ over $A_k$. The set $A_k$ is a copy of $B$ or of $C$; it is closed in $T_k$ and in
$T_{k-1}$, by hypothesis and, for a piece that is a copy of $E$, by
Lemma~\ref{lem:calc}(c). By induction on $k$ and Lemma~\ref{lem:calc}(c),(d),
$L_k$ is positive and each of $T_0,\dots,T_k$ is closed in $L_k$. So $L$ is
positive and every $T_k$ is closed in $L$. By Lemma~\ref{lem:bridgeset}(iv) the
map $q$ identifies at most $m$ pairs of points, and every hyperedge of $Z$ is
the image of a hyperedge of $L$. Hence, for $U\subseteq Z$ and $\widetilde U=q^{-1}(U)$,
\begin{equation}\label{eq:cutloss}
 \delta_Z(U)\ \ge\ \delta_L(\widetilde U)-2m .
\end{equation}

\emph{Positivity and closedness.} Let $U\subseteq Z$ be non-empty.

\emph{Case 1: some piece $Z_r$ contains no point of $U$ other than
$F$-points.} Cut between $r-1$ and $r$. No point of $U$ is then the image of
two points of $L$, so $q$ maps $\widetilde U$ bijectively onto $U$ and
$E_L(\widetilde U)$ bijectively onto $E_Z(U)$. Hence
$\delta_Z(U)=\delta_L(\widetilde U)>0$. If $U$ properly contains a piece $Z_i$,
then $\widetilde U$ properly contains the corresponding $T_k$, and
$\delta_Z(U)=\delta_L(\widetilde U)>\delta_L(T_k)=\delta_Z(Z_i)$.

\emph{Case 2: every piece contains a point of $U$ that is not an $F$-point.} Cut
between $s-1$ and $0$, and put $\widetilde U_k=\widetilde U\cap T_k$. Applying
\eqref{eq:amalgam} to $L_k=L_{k-1}\cup T_k$ for $k=1,\dots,s-1$ gives
\[
 \delta_L(\widetilde U)=\delta(\widetilde U_0)
 +\sum_{k=1}^{s-1}\bigl(\delta(\widetilde U_k)-\delta(\widetilde U_k\cap A_k)\bigr).
\]
By Lemma~\ref{lem:calc}(b), $\widetilde U_k\cap A_k\cle\widetilde U_k$, so the
$k$-th summand is non-negative, and it is at least $1$ when
$\widetilde U_k\not\subseteq A_k$; also $\delta(\widetilde U_0)\ge 1$ when
$\widetilde U_0\ne\varnothing$. Let $J$ be the set of those $k$ for which
$\widetilde U_k$ contains a point that is not an $F$-point and does not lie in
$T_{k-1}$ (no condition on $T_{k-1}$ for $k=0$). Each $k\in J$ contributes at
least $1$, so $\delta_L(\widetilde U)\ge|J|$. Let $0\le k'<s$. By assumption
$T_{k'}$ contains a point $u\in\widetilde U$ that is not an $F$-point; the set
of $k$ with $u\in T_k$ consists of at most three consecutive values, and its
least element $k$ belongs to $J$, with $k\le k'\le k+2$. Hence $3|J|\ge s$, and
by \eqref{eq:cutloss} and \eqref{eq:Nbound}
\[
 \delta_Z(U)\ \ge\ \frac s3-2m\ >\ \delta(E)\ \ge\ \delta(D)\ >\ 0 ,
\]
where $\delta(E)\ge\delta(D)$ because $d_0(D)\cle E$. In particular
$\delta_Z(U)>\delta_Z(Z_i)$ for every piece $Z_i$.

In both cases $\delta_Z(U)>0$, and $\delta_Z(U)>\delta_Z(Z_i)$ whenever $U$
properly contains the piece $Z_i$. So $Z$ is positive and every piece is closed
in $Z$; with the first step, $e$ and all $a_i$ are closed embeddings.
\end{proof}

\begin{proposition}\label{prop:stabdelta}
For all $B,C\in\K$ we have $\langle G_B,G_C\rangle=G_{B\cap C}$.
\end{proposition}

\begin{proof}
One inclusion is clear. Let $g\in G_{B\cap C}$. Choose $D\in\K$ containing
$B\cup C$ and $E\in\K$ containing $D\cup gD$. Then $B\cle D$ and $C\cle D$ by
Lemma~\ref{lem:calc}(b), and the inclusion $d_0\colon D\to E$ and
$d_1=g|_D\colon D\to E$ are closed embeddings that agree on $B\cap C$. Let
$Z$, $e$ and $a_0,\dots,a_{2N}$ be as in Lemma~\ref{lem:bridgedelta}. By (M2)
there is an embedding $h\colon Z\to M$ with closed image and $he=\id_E$. Each
$ha_i\colon D\to M$ is an isomorphism onto a closed subset, so by (M3)
$ha_i=g_i|_D$ for some $g_i\in G$; since $ha_0=d_0$ and $ha_{2N}=d_1$ we may take
$g_0=1$ and $g_{2N}=g$. The relations between consecutive $a_i$ say that
$g_i^{-1}g_{i+1}$ lies in $G_B$ for $i$ even and in $G_C$ for $i$ odd. Hence
$g=\prod_{i=0}^{2N-1}g_i^{-1}g_{i+1}\in\langle G_B,G_C\rangle$.
\end{proof}

\begin{corollary}\label{cor:opendelta}
$(M,G,\K)$ satisfies \textup{(K3)}: every open subgroup of $G$ is $G_F$ for some
$F\in\K$.
\end{corollary}

\begin{proof}
Let $U$ be an open subgroup. By (K1) it contains $G_B$ for some $B\in\K$;
choose such a $B$ of least cardinality. For $g\in U$ we have
$G_{gB}=gG_Bg^{-1}\le U$, so $G_{B\cap gB}\le U$ by
Proposition~\ref{prop:stabdelta}; here $B\cap gB\in\K$ by
Lemma~\ref{lem:calc}(b),(c). By minimality $gB=B$, and an increasing bijection
of a finite linearly ordered set is the identity, so $g\in G_B$. Hence $U=G_B$.
\end{proof}

\subsection{The two properties}

\begin{theorem}\label{thm:total}
The category $\Cd$ is total: it is non-empty and connected, and for every
functor $L\colon\Cd\to\Set$ every connected component of $\El(L)$ is weakly
contractible.
\end{theorem}

\begin{proof}
By Lemmas~\ref{lem:AMequiv}, \ref{lem:K2delta}, \ref{lem:SCdelta} and
Corollary~\ref{cor:opendelta} the triple $(M,G,\K)$ satisfies (K1)--(K4). By
Proposition~\ref{prop:reduction} every action component of $\mathcal A$ is
weakly contractible or weakly equivalent to $\mathcal A$, and $B\mathcal A$ is
weakly contractible by Lemma~\ref{lem:SCdelta}. Since $\mathcal A$ is equivalent
to $\Cd$, the same holds for $\Cd$, and Proposition~\ref{prop:total} applies.
\end{proof}

\begin{theorem}\label{thm:notexact}
Let $R$ be a non-zero ring. For an object $X$ of $\Cd$ let
$P_X=R[\Cd(X,-)]$, a functor from $\Cd$ to left $R$-modules, and let
\[
 \Theta\colon P_W\longrightarrow P_V\oplus P_V,\qquad
 \Theta(g)=(gv_2-gv_1,\ gv_3-gv_1)\quad(g\colon W\to X).
\]
Then $\Theta$ is a monomorphism, and $\colim\Theta\colon R\to R\oplus R$ is zero.
In particular the functor $\colim$ from $\Cd$-indexed diagrams of $R$-modules to
$R$-modules is not exact; for $R=\Z$, the category $\Cd$ is not exact.
\end{theorem}

\begin{proof}
Let $X$ be an object and let $t=\sum_g c_g\,g$ be a non-zero element of
$P_W(X)$; let $g_1,\dots,g_r$ be the arrows $W\to X$ with $c_g\ne 0$. An
increasing map of $W$ is determined by its image, so the images of the $g_k$
are $r$ distinct hyperedges of $X$. Let $U$ be their union and $v=|U|$. Then
$0<\delta(U)\le 2v-3r$. The $r$ hyperedges have $3r<2v$ incidences with the $v$
points of $U$, each of which lies on at least one of them; so some point $u$
lies on exactly one, say on the image of $g_k$, where $u=g_k(i)$. The arrow
$g_kv_i\colon V\to X$ has image $\{u\}$, and no other arrow $g_lv_j$ with
$c_{g_l}\ne 0$ has this image. The coefficient of $g_kv_i$ is therefore
$c_{g_k}$ in the first component of $\Theta(t)$ if $i=2$, it is $c_{g_k}$ in the
second component if $i=3$, and it is $-c_{g_k}$ in both if $i=1$. So
$\Theta(t)\ne 0$.

The colimit of a representable functor $P_X$ is $R$, the canonical map
$P_X(Y)\to R$ being the augmentation $\sum c_f f\mapsto\sum c_f$. Both
components of $\Theta(g)$ have augmentation zero, so $\colim\Theta=0$, which is
not injective because $\colim P_W=R\ne 0$.
\end{proof}
% ---- end sec4_cdelta.tex ----
% ---- begin sec5_twographs.tex (inlined for the deposit) ----
\section{A quasi-filtered category that is not total}\label{sec:twographs}

\subsection{Ordered two-graphs}

Let $S$ be a finite set. A \emph{graph} on $S$ is a function $b$ from the
two-element subsets of $S$ to $\F_2$. A \emph{two-graph} on $S$ is a function
$c$ from the three-element subsets of $S$ to $\F_2$ such that, for every
four-element subset $Q$ of $S$, the sum of $c$ over the four three-element
subsets of $Q$ is zero. (The triples with $c=1$ then form a two-graph in the
usual sense; see \cite{Seidel}.) We write $c(x,y,z)$ for $c(\{x,y,z\})$ and
put $c(x,y,z)=0$ when two of the entries coincide; similarly for graphs. For
a graph $b$ let
\[
 (\partial b)(x,y,z)=b(x,y)+b(x,z)+b(y,z).
\]
This is a two-graph, because in the sum over the four triples of a
four-element set every pair occurs twice. For a two-graph $c$ on $S$ and
$p\in S$ let $c_p$ be the graph $c_p(x,y)=c(p,x,y)$. Then
\begin{equation}\label{eq:anchor}
 \partial(c_p)=c ;
\end{equation}
for a triple that does not contain $p$ this is the condition on the
four-element set formed with $p$, and for a triple $\{p,y,z\}$ it holds because
$c_p(p,y)=c_p(p,z)=0$.

An \emph{ordered two-graph} is a finite non-empty linearly ordered set with a
two-graph. An \emph{embedding} of ordered two-graphs is an increasing
injection $f\colon X\to Y$ with $c_Y(f(x),f(y),f(z))=c_X(x,y,z)$ for all
$x,y,z$.

\begin{definition}
The objects of $\mathcal T$ are the ordered two-graphs with underlying ordered
set $[n]$, $n\ge1$; its arrows are the embeddings.
\end{definition}

\begin{lemma}[anchored amalgams]\label{lem:anchored}
Let $S=S_1\cup S_2$ and let $c_i$ be a two-graph on $S_i$, $i=1,2$, such that
$c_1$ and $c_2$ agree on the triples contained in $S_1\cap S_2$. Choose
$p_i\in S_i$, with $p_1=p_2\in S_1\cap S_2$ if this set is non-empty, and let
$b$ be the graph on $S$ that equals $(c_i)_{p_i}$ on the pairs contained in
$S_i$ and is zero on all other pairs. Then $\partial b$ is a two-graph on $S$
that restricts to $c_i$ on $S_i$.
\end{lemma}

\begin{proof}
The graph $b$ is well defined, because $(c_1)_p$ and $(c_2)_p$ agree on the
pairs contained in $S_1\cap S_2$ when $p$ lies in this set. On triples
contained in $S_i$ we have $\partial b=\partial((c_i)_{p_i})=c_i$ by
\eqref{eq:anchor}.
\end{proof}

We call $\partial b$ the amalgam \emph{anchored at} $p_1,p_2$. In particular
any two objects of $\mathcal T$ embed into a third, so $\mathcal T$ is
connected.

\subsection{The generic ordered two-graph}

\begin{proposition}\label{prop:generic2}
There is an ordered two-graph structure on a countably infinite set $M$
\textup{(}a linear order and a function $c$ on three-element subsets that is a
two-graph on every finite subset\textup{)} such that:
\begin{enumerate}
\item[\textup{(N2)}] if $Y\subseteq M$ is finite and $f\colon Y\to C$ is an
 embedding into an ordered two-graph, there is an embedding $h\colon C\to M$
 with $hf=\id_Y$ \textup{(}here $Y$ may be empty\textup{)};
\item[\textup{(N3)}] every isomorphism between finite subsets of $M$ extends to an
 automorphism of $M$.
\end{enumerate}
\end{proposition}

\begin{proof}
As for Proposition~\ref{prop:generic}, with every finite subset in the role
of a closed subset and with the anchored amalgams of Lemma~\ref{lem:anchored},
ordered by any linear order extending the two given ones, in the role of the
free amalgams.
\end{proof}

Fix such an $M$, let $G=\Aut(M)$ and let $\K$ be the set of all finite subsets
of $M$. By (N3) the arrows of $\mathcal A=\mathcal A(M,G,\K)$ are the embeddings
between finite non-empty subsets of $M$, and by (N2) with $Y=\varnothing$ the
category $\mathcal A$ is equivalent to $\mathcal T$. Condition (K1) is clear.
By (N2) the order of $M$ has no greatest element.

\begin{lemma}\label{lem:K24two}
$(M,G,\K)$ satisfies \textup{(K2)} and \textup{(K4)}.
\end{lemma}

\begin{proof}
(K2). Let $B\subseteq M$ be finite and $x\in M\setminus B$. Let
$C=B\cup\{x,x'\}$ carry an amalgam of $B\cup\{x\}$ and of a second copy
$B\cup\{x'\}$ of it over $B$ (Lemma~\ref{lem:anchored}), ordered with $x'$
immediately after $x$. By (N2) it embeds into $M$ over $B\cup\{x\}$, and by
(N3) the isomorphism $B\cup\{x\}\to B\cup\{x'\}$ extends to some $g\in G_B$
with $gx\ne x$.

(K4). We argue as for Lemma~\ref{lem:SCdelta}. Let $\mathcal C$ be the
category of ordered two-graphs whose underlying set is a finite subset of
$\Omega$, with the embeddings as arrows; it is equivalent to $\mathcal A$. Let
$P$ be an object, $p_0$ its least element, $\mathcal U$ a non-empty cosieve of
$P\backslash\mathcal C$ and $(Z,z)$ an object of $\mathcal U$. For an object
$(X,x)$ let $T(X)$ be the set $\lambda(X)\cup\rho(Z\setminus z(P))$, ordered
as in Lemma~\ref{lem:SCdelta}, with the amalgam of the copies of $X$ and of
$Z$ anchored at $\lambda(x(p_0))$. For an arrow $f\colon X\to X'$ under $P$,
the map $T(f)$ defined as before is increasing. It carries the anchor to the
anchor and the pairs contained in one of the two copies to such pairs,
preserving the values of the graph $b$ of Lemma~\ref{lem:anchored}, and it
carries the remaining pairs, on which $b=0$, to such pairs. Hence it
preserves $\partial b$ and is an embedding. As before we obtain an
endofunctor $T$ of $P\backslash\mathcal C$ and natural transformations
$\id\to T\leftarrow\mathrm{const}_{(Z,z)}$, and $B\mathcal U$ is contractible.
\end{proof}

\subsection{Open subgroups}

\begin{lemma}\label{lem:stab2a}
Let $B,C\subseteq M$ be finite with $B\cap C\ne\varnothing$. Then
$\langle G_B,G_C\rangle=G_{B\cap C}$.
\end{lemma}

\begin{proof}
Let $g\in G_{B\cap C}$. Put $D=B\cup C$, $E=D\cup gD$, $F=B\cap C$, and let
$d_0\colon D\to E$ be the inclusion and $d_1=g|_D$. Let $Z$, $e$,
$a_0,\dots,a_s$ and the pieces $Z_i$ be as in Lemma~\ref{lem:bridgeset}. Each
piece is a copy of $E$ or of $D$ and carries the corresponding two-graph.
Two consecutive pieces induce the same two-graph on their intersection, which
is the common image of $B$ or of $C$, because $d_0$ and $d_1$ are embeddings;
call this $(\ast)$. Fix $p\in F$ and let $\hat p$ be the corresponding
$F$-point of $Z$. Let $b$ be the graph on $Z$ with
$b(u,u')=c_{Z_i}(\hat p,u,u')$ if $u$ and $u'$ lie in a common piece $Z_i$, and
$b(u,u')=0$ if they lie in no common piece. This is well defined. Indeed, let
$u,u'\in Z_i\cap Z_j$. If $i$ and $j$ are consecutive, apply $(\ast)$. If not,
and $u$, $u'$ are both $F$-points, note that all pieces induce the same
two-graph on the set of $F$-points. Otherwise $\{i,j\}=\{1,s-1\}$ and
$u,u'\in Z_0$ by Lemma~\ref{lem:bridgeset}(ii), and we apply $(\ast)$ to
$(Z_j,Z_0)$ and to $(Z_0,Z_i)$. Give $Z$ the two-graph $\partial b$. On a
triple contained in a piece $Z_i$ it equals $\partial((c_{Z_i})_{\hat p})=c_{Z_i}$
by \eqref{eq:anchor}. So $e$ and all $a_i$ are embeddings of ordered
two-graphs.

By (N2) there is an embedding $h\colon Z\to M$ with $he=\id_E$, and by (N3)
$ha_i=g_i|_D$ with $g_i\in G$, $g_0=1$, $g_s=g$. By
Lemma~\ref{lem:bridgeset}(i), $g_i^{-1}g_{i+1}$ lies in $G_B$ for $i$ even and
in $G_C$ for $i$ odd, so $g\in\langle G_B,G_C\rangle$.
\end{proof}

\begin{lemma}\label{lem:stab2b}
For distinct points $x,y$ of $M$ we have $\langle G_{\{x\}},G_{\{y\}}\rangle=G$.
\end{lemma}

\begin{proof}
Let $H$ be the subgroup on the left, and suppose $x<y$. A two-graph on a set
with at most two elements is empty, so by (N3) the group $G$ is transitive on
$M$, and $G_{\{w\}}$ is transitive on $\{u:u<w\}$ and on $\{u:u>w\}$ for every
$w\in M$. Note that $H$ contains $G_{\{hx\}}$ and $G_{\{hy\}}$ for every
$h\in H$. Let $u\in M$ and choose $w$ greater than $u$ and $y$. Then
$w\in G_{\{x\}}\,y$, so $G_{\{w\}}\le H$, and $u\in G_{\{w\}}\,x$. Hence $H$
is transitive on $M$. For $g\in G$ choose $h\in H$ with $hx=gx$; then
$h^{-1}g\in G_{\{x\}}\le H$, so $g\in H$.
\end{proof}

\begin{proposition}\label{prop:open2}
For all finite $B,C\subseteq M$ we have $\langle G_B,G_C\rangle=G_{B\cap C}$, and
$(M,G,\K)$ satisfies \textup{(K3)}.
\end{proposition}

\begin{proof}
If $B\cap C\ne\varnothing$ this is Lemma~\ref{lem:stab2a}. Otherwise choose
distinct points $x,y$ outside $B\cup C$. By Lemma~\ref{lem:stab2a},
$\langle G_B,G_C\rangle$ contains
$\langle G_{B\cup\{x\}},G_{C\cup\{x\}}\rangle=G_{\{x\}}$, and likewise
$G_{\{y\}}$; so it is $G$ by Lemma~\ref{lem:stab2b}. Condition (K3) follows as
in Corollary~\ref{cor:opendelta}.
\end{proof}

\subsection{The classifying space}

\begin{lemma}\label{lem:pi1}
$B\mathcal T$ is simply connected.
\end{lemma}

\begin{proof}
$\mathcal T$ is connected, so by the correspondence between covering spaces
and functors that invert all arrows (Section~\ref{sec:classes}) it
suffices to show that every functor $\Psi\colon\mathcal T\to\Set$ that sends
all arrows to bijections is isomorphic to a constant functor. Let $P$ and
$Q$ be the objects with one and with two points (there are no others with at
most two points), and $d^1,d^0\colon P\to Q$ the arrows onto the first and the
second point. Let $X$ be an object with three points and
$e_{ij}\colon Q\to X$ the embedding onto $\{i,j\}$. Then
\[
 e_{12}d^1=e_{13}d^1,\qquad e_{13}d^0=e_{23}d^0,\qquad e_{12}d^0=e_{23}d^1 .
\]
Since $\Psi$ takes values in bijections, the first two equations give
$\Psi(e_{12})=\Psi(e_{13})=\Psi(e_{23})$, and then the third gives
$\Psi(d^0)=\Psi(d^1)$. For an object $Y$ and $u\in Y$ let $i_u\colon P\to Y$ be
the arrow with image $\{u\}$. If $u<u'$ and $\epsilon\colon Q\to Y$ is the
embedding onto $\{u,u'\}$, then $i_u=\epsilon d^1$ and $i_{u'}=\epsilon d^0$, so
$\Psi(i_u)=\Psi(i_{u'})$. Thus $\theta_Y=\Psi(i_u)\colon\Psi(P)\to\Psi(Y)$ does
not depend on $u$. For an arrow $f\colon Y\to Y'$ we have $fi_u=i_{f(u)}$,
hence $\Psi(f)\theta_Y=\theta_{Y'}$. So $\theta$ is a natural isomorphism from
the constant functor with value $\Psi(P)$ to $\Psi$.
\end{proof}

\begin{lemma}\label{lem:H2}
$H_2(B\mathcal T;\F_2)\ne 0$.
\end{lemma}

\begin{proof}
We use the chains of the nerve with coefficients in $\F_2$. For a
$2$-simplex $\sigma=(X_0\xrightarrow{f_1}X_1\xrightarrow{f_2}X_2)$ let
$m_0\le m_1\le m_2$ be the images in $X_2$ of the greatest elements of $X_0$,
$X_1$, $X_2$, and put $\Phi(\sigma)=c_{X_2}(m_0,m_1,m_2)$. Then $\Phi$ is a
cocycle. Indeed, for a $3$-simplex $X_0\to X_1\to X_2\to X_3$ let
$n_0\le n_1\le n_2\le n_3$ be the images in $X_3$ of the greatest elements.
Since $X_2\to X_3$ is an embedding, the values of $\Phi$ on the four faces are
the values of $c_{X_3}$ on the four triples obtained by omitting one of the
$n_k$. If the $n_k$ are distinct, their sum is zero by the definition of a
two-graph. If $n_k=n_{k+1}$ for some $k$, the triples obtained by omitting
$n_k$ or $n_{k+1}$ give the same value, and the other two contain a
repetition.

Let $X$ and $X'$ be the two objects with three points, with
$c_X(1,2,3)=1$ and $c_{X'}(1,2,3)=0$. For $S_0\subset S_1\subset[3]$ with
$|S_0|=1$, $|S_1|=2$ let $\sigma(S_0,S_1,X)$ be the $2$-simplex
$P\to Q\to X$ in which the image of $Q$ is $S_1$ and the image of $P$ is
$S_0$; define $\sigma(S_0,S_1,X')$ likewise. Let $\xi$ be the sum of these
twelve simplices. It is a cycle: the face $P\to Q$ does not depend on the
target $X$ or $X'$; the face $Q\to X$ depends only on $S_1$ and so occurs
twice, and the face $P\to X$ depends only on $S_0$ and occurs twice; similarly
for $X'$. By definition $\Phi(\sigma(S_0,S_1,X))=c_X(\max S_0,\max S_1,3)$,
which is non-zero only for $S_0=\{1\}$, $S_1=\{1,2\}$, and $\Phi$ vanishes on
the simplices with target $X'$. So $\Phi(\xi)=1$, and $\xi$ is not a boundary.
\end{proof}

\begin{theorem}\label{thm:twographs}
Every action component of $\mathcal T$ is weakly contractible or weakly
equivalent to $\mathcal T$. Consequently $\mathcal T$ is quasi-filtered, and
it is not total.
\end{theorem}

\begin{proof}
By Lemma~\ref{lem:K24two} and Proposition~\ref{prop:open2} the triple
$(M,G,\K)$ satisfies (K1)--(K4), so the first assertion holds for $\mathcal A$
by Proposition~\ref{prop:reduction}, and hence for the equivalent category
$\mathcal T$. By Lemma~\ref{lem:pi1} all action components of $\mathcal T$ are
simply connected (each is weakly contractible or weakly equivalent to
$\mathcal T$), so $\mathcal T$ is quasi-filtered by
Proposition~\ref{prop:qf}. The category of elements of the functor with value
a point is $\mathcal T$ itself, which is not weakly contractible by
Lemma~\ref{lem:H2}; so $\mathcal T$ is not total by
Proposition~\ref{prop:total}.
\end{proof}

\begin{remark}\label{rem:tower}
The run from which this example is taken identifies $B\mathcal T$ with an
Eilenberg--Mac~Lane space $K(\F_2,2)$, and treats $n$-cocycles in the same way
for every $n\ge2$; see Section~\ref{sec:provenance}. Those statements are not
used and not claimed here.
\end{remark}
% ---- end sec5_twographs.tex ----
% ---- begin sec6_manuscript.tex (inlined for the deposit) ----
\section{The passage of \emph{Les D\'erivateurs}}\label{sec:manuscript}

In Chapter~II of \cite{Der} (manuscript pages 34--36ter; the last of these is
on the printed page 23, page 25 of the PDF, of the edition's file for that
chapter), Grothendieck shows that filtered categories are \emph{totalement
$W$-asph\'eriques} (Corollaire~2.17; for a functor $f\colon X\to Y$ the
condition is on the category of elements of the presheaf
$y\mapsto\pi_0(y\backslash X)$, which is the condition of
Definition~\ref{def:classes}(c)). He then asks for the converse, for the
class $W_\infty$ of weak equivalences: ``cela implique-t-il que $X$ soit
filtrante~? \c{C}a ne m'\'etonnerait pas'' (ms.~p.~36). His Proposition~2.19
answers in the negative: it asserts that $X$ is totally $W_\infty$-aspherical
whenever the diagonal $X\to X\times X$ is final (\emph{cofinal}, in his
terminology), that is, whenever $X$ is sifted; and a corollary states that
the opposite $\boldsymbol\Delta^{\circ}$ of the category of simplices is
totally $W_\infty$-aspherical and not filtered. In his Contre-exemple~2.20
(ms.~pp.~36bis--36ter) he derives from Proposition~2.19 that every connected
simplicial set would be weakly contractible, and concludes: ``Cela me fait
soup\c{c}onner que la conjecture
$X$ totalement $W_\infty$-asph\'erique $\Longrightarrow$ $X$ filtrante
pourrait \^etre valable.''

The results above give the following answers.
\begin{enumerate}
\item The conjecture is false. The category $\Delta_{\mathrm{inj}}$ is exact
 by the theorems of Isbell and Mitchell \cite[Theorems~1 and~3]{IM73}, hence total
 by Proposition~\ref{prop:incl}, and it is not filtered. The category $\Cd$
 is a second example, which is not even exact.
\item In the edition, Proposition~2.19 is marked false (footnote~43, and
 footnote~45 at the step of its proof that fails), and Grothendieck wrote his
 Contre-exemple~2.20 ``\`a la proposition fausse 2.19'' against it; the corollary falls with it, and indeed
 $\boldsymbol\Delta^{\circ}$ is not total. This follows from
 Proposition~\ref{prop:total}, since the action components of
 $\boldsymbol\Delta^{\circ}$ are the opposites of the categories of simplices
 of connected simplicial sets; it also follows from
 Proposition~\ref{prop:sifted}, since $\boldsymbol\Delta^{\circ}$ is sifted
 and not filtered. More generally, by Proposition~\ref{prop:sifted} a sifted
 category is total only if it is filtered. (His Contre-exemple~2.20 uses
 Chapter~XII, Th\'eor\`eme~4, Corollaire~4, which does not hold for
 $W_\infty$: where it concludes that a category is weakly contractible, only
 the vanishing of its reduced integral homology follows. His conclusion
 stands, and the two arguments above do not use that step.)
\item The conjecture holds for sifted categories
 (Proposition~\ref{prop:sifted}), and it holds for categories with a weakly
 terminal object: there even quasi-filtered categories are filtered, by a
 theorem of Isbell and Mitchell \cite[Corollary~4.4]{IM76}.
\end{enumerate}
% ---- end sec6_manuscript.tex ----
% ---- begin sec7_provenance.tex (inlined for the deposit) ----
\section{Provenance and checks}\label{sec:provenance}

This paper is AI-produced. No human has proved, written or checked any part
of it.

\paragraph{Sources.} The mathematics was found in four runs of GPT-6 Astra
(OpenAI, ChatGPT Pro mode), in November 2026. The run TOTAL\_ASPHERICITY proved
Proposition~\ref{prop:total} (its statement C.5) and
Proposition~\ref{prop:incl}. The run PRESHEAF\_TOTALITY proved
Proposition~\ref{prop:sifted} (its statement F.3). The run TOTAL\_EXACT\_2
found the category $\Cd$ and proved Theorem~\ref{thm:A} (its statements A.1,
A.3, with the proofs A.21, A.24, A.30--A.34), and stated
Theorem~\ref{thm:C} (D.1, D.2) with another third example, a modification of
a construction of \cite{AKR}, whose proof is long and is not reproduced or
used here. The run DER\_FOLLOWUP found the cocycle categories, of which
$\mathcal T$ is the first, and proved Theorem~\ref{thm:B} (its statements
A.3--A.8), for cocycles of every degree $n\ge 2$.

\paragraph{The text.} The present text was written by Claude Fable~5.2 at
effort max. The proofs of Sections~\ref{sec:cdelta} and~\ref{sec:reduction}
follow those of the runs, reorganized: the reduction is stated once
(Proposition~\ref{prop:reduction}) and the bridge is split into a statement
on ordered sets (Lemma~\ref{lem:bridgeset}) and its two uses. The following
arguments are the writer's and differ from the runs': the proofs of
Propositions~\ref{prop:qf} and~\ref{prop:sifted} in the form given; the
reduction of Proposition~\ref{prop:open2} to the case of a non-empty
intersection, with Lemmas~\ref{lem:stab2a} and~\ref{lem:stab2b}; and
Lemmas~\ref{lem:pi1} and~\ref{lem:H2}.

\paragraph{Checks.} The proofs of the run TOTAL\_EXACT\_2 used here were
followed in full by two readers before this text was written (Claude
Fable~5.2 at effort high, and Claude Fable~5.2 at effort max, 13 November
2026), and those of DER\_FOLLOWUP by one (Claude Fable~5.2 at effort max, 17
November 2026). Every proof in this text was derived again by the writer
while writing. A program independent of the runs searched all objects of
$\Cd$ with at most seven vertices for an integral, rational or mod-$p$
($p=2,3,5$) affine equalizer of $v_1,v_2,v_3$ and found none, as
Theorem~\ref{thm:notexact} predicts. The finished text was then checked
independently by Claude Opus~5.5 at effort max (25 November 2026): every
statement and every proof was found correct, the fact that $\mathcal T$ is not
total was proved again by another route, and eight repairs of wording and
citation were made, none of which changes a statement. So the proofs taken
from the runs have been checked by models of a family other than the one
that found them, and the arguments that are the writer's own have been
checked by one other model, of the writer's family. Nothing is formalized.
The record of every check, with model, effort and date, is in the file
\texttt{CHECK\_EVENTS.tsv} that accompanies this paper in its workbench, reached
from \url{https://github.com/KokunoYumeto/mathematics-workbenches};
\texttt{PROVENANCE.md} lists every contributor and says in which copy each
citation was confirmed, and which citations are given only at the level of
a table of contents.
% ---- end sec7_provenance.tex ----

\begin{thebibliography}{99}
\bibitem{AKR} J.~Ad\'amek, V.~Koubek and J.~Reiterman, Embeddings into
 categories with fixed points in representations, Czechoslovak Math.\ J.\
 31(106) (1981), no.~3, 368--389.
\bibitem{AR} J.~Ad\'amek and J.~Rosick\'y, On sifted colimits and generalized
 varieties, Theory Appl.\ Categ.\ 8 (2001), no.~3, 33--53.
\bibitem{BK} A.~K.~Bousfield and D.~M.~Kan, \emph{Homotopy limits,
 completions and localizations}, Lecture Notes in Math.\ 304, Springer, 1972.
\bibitem{EHN} D.~M.~Evans, J.~Hubi\v{c}ka and J.~Ne\v{s}et\v{r}il,
 Automorphism groups and Ramsey properties of sparse graphs, Proc.\ London
 Math.\ Soc.\ 119 (2019), 515--546; arXiv:1801.01165.
\bibitem{GJ} P.~G.~Goerss and J.~F.~Jardine, \emph{Simplicial homotopy
 theory}, Progress in Math.\ 174, Birkh\"auser, 1999.
\bibitem{Der} A.~Grothendieck, \emph{Les D\'erivateurs}, manuscript,
 1990--1991; edited by M.~K\"unzer, J.~Malgoire and G.~Maltsiniotis,
 \url{https://webusers.imj-prg.fr/~georges.maltsiniotis/groth/Derivateurs.html}.
\bibitem{IM73} J.~Isbell and B.~Mitchell, Exact colimits, Bull.\ Amer.\
 Math.\ Soc.\ 79 (1973), 994--996.
\bibitem{IM76} J.~Isbell and B.~Mitchell, Exact colimits and fixed points,
 Trans.\ Amer.\ Math.\ Soc.\ 220 (1976), 289--298.
\bibitem{JardineL8} J.~F.~Jardine, Lecture 08: Bisimplicial sets, homotopy
 limits and colimits, lecture notes on homotopy theory, on the author's page,
 \url{https://math.sci.uwo.ca/~jardine/HomTh/lecture008a.pdf}.
\bibitem{CWM} S.~Mac~Lane, \emph{Categories for the working mathematician},
 second edition, Graduate Texts in Math.\ 5, Springer, 1998.
\bibitem{Quillen} D.~Quillen, Higher algebraic $K$-theory~I, in
 \emph{Algebraic $K$-theory~I}, Lecture Notes in Math.\ 341, Springer, 1973,
 85--147.
\bibitem{Rosicky} J.~Rosick\'y, On homotopy varieties, Adv.\ Math.\ 214
 (2007), 525--550; Corrigendum, Adv.\ Math.\ 259 (2014), 841--842;
 arXiv:math/0509655v2 (the corrected version, whose numbering is used here).
\bibitem{Seidel} J.~J.~Seidel, A survey of two-graphs, in \emph{Colloquio
 Internazionale sulle Teorie Combinatorie (Roma, 1973)}, Tomo~I, Atti dei
 Convegni Lincei 17, Accad.\ Naz.\ Lincei, 1976, 481--511.
\bibitem{Weibel} C.~A.~Weibel, \emph{An introduction to homological
 algebra}, Cambridge Studies in Advanced Math.\ 38, Cambridge University
 Press, 1994.
\bibitem{Kbook} C.~A.~Weibel, \emph{The $K$-book: an introduction to
 algebraic $K$-theory}, Graduate Studies in Math.\ 145, Amer.\ Math.\ Soc.,
 2013; Chapter~IV is cited from the version on the author's page,
 \url{https://sites.math.rutgers.edu/~weibel/Kbook/Kbook.IV.pdf}.
\end{thebibliography}
\end{document}
